% LaTeX template: mahmoud.s.fahmy@students.kasralainy.edu.eg
% For more details: https://www.sharelatex.com/learn/Beamer
\documentclass[10pt]{beamer}					% Document class
\usepackage[english]{babel}				% Set language
\usepackage[utf8x]{inputenc}			% Set encoding
\usepackage{color}
\usepackage[normalem]{ulem}
\mode<presentation>						% Set options
{
  \usetheme{default}					% Set theme
  \usecolortheme{default} 				% Set colors
  \usefonttheme{default}  				% Set font theme
  \setbeamertemplate{caption}[numbered]	% Set caption to be numbered
}
% Uncomment this to highlight the current section in the outline at the start of each section.
%\AtBeginSection[]
%{
%  \begin{frame}{Outline}
%    \tableofcontents[currentsection]
%  \end{frame}
%}
\usepackage{graphicx}					% For including figures
\usepackage{url}
\usepackage{booktabs}					% For table rules
\usepackage{hyperref}					% For cross-referencing
\input{macro}

\title{Progress on the ADAPT RCT}	% Presentation title
\author{Xueqing Liu}								% Presentation author
% \institute{Harvard University}					% Author affiliation
\date{\today}									% Today's date
\begin{document}

\maketitle

\begin{frame}{Outline}
    \tableofcontents
\end{frame}

\section{Previously\ldots}

\begin{frame}{Previously\ldots}
On the vanilla testbed, all RL variants achieve similar performance. We attribute this to two factors:
    \begin{itemize}
        \item the standardized treatment effect (STE) is too small on the vanilla testbed;
        \item the RL algorithm carries too many state variables.
    \end{itemize}
\end{frame}

\section{Mediators}

\begin{frame}{Mediators related to perceived utility $E_w$}
    \begin{enumerate}
        \item Unprompted app page views over the next 4 hours, $M^{E,PV}_{w,d,t} \in \mathbb{Z}$.
        \item Next-morning Fitbit wear status, $M^{E,FW}_{w,d+1}\in \{0,1\}$.
        \begin{itemize}
            \item Measured over the 2 hours preceding the morning decision time.
            \item Defined as ``worn'' if heart rate is positive or step count is positive for at least 100 minutes.
            \item \textbf{Issue:} the new Fitbit (Google Health) omits rows in which heart rate is 0, so we cannot distinguish ``not worn'' from ``not synced.''
        \end{itemize}
        \item Next-morning check-in completion status, $M^{E,PJ}_{w,d+1} \in \{0,1\}$.
        \begin{itemize}
            \item The morning check-in window is 6\,am to 11\,am.
            \item $M^{E,PJ}_{w,d+1}$ can be observed only between the morning and afternoon decision times.
            \item When used in the state for $A_{w,d+1,1}$ and $A_{w,d+1,2}$, we split it into $M^{E,PJ}_{w,d+1,1}$ and $M^{E,PJ}_{w,d+1,2}$.
            \item \textbf{Question:} should we treat $M^{E,PJ}_{w,d,1}$ and $M^{E,PJ}_{w,d,2}$ as the mediators, where $M^{E,PJ}_{w,d,1}$ mediates the previous day's two actions $A_{w,d-1,1}$ and $A_{w,d-1,2}$, while $M^{E,PJ}_{w,d,2}$ mediates both the previous day's actions and this morning's action $A_{w,d,1}$?
        \end{itemize}
    \end{enumerate}
\end{frame}

\begin{frame}{Mediators related to CAE $Y_w$}
    \begin{enumerate}
        \item Active minutes over the next 4 hours, $M^{Y,SC}_{w,d,t} \in \mathbb{Z}$.
        \item Anticipated affect from the next-morning check-in, $M^{Y,AA}_{w,d+1}$.
        \begin{itemize}
            \item As with $M^{E,PJ}_{w,d+1}$, it can be observed only between the morning and afternoon decision times.
            \item When used in the state for $A_{w,d+1,1}$ and $A_{w,d+1,2}$, we split it into
            $M^{Y,AA}_{w,d+1,1}\cdot M^{E,PJ}_{w,d+1} \cdot I_{w,d+1}^{\text{day}}$ and
            $M^{Y,AA}_{w,d+1,2}\cdot M^{E,PJ}_{w,d+1} \cdot I_{w,d+1}^{\text{day}}$,
            where $I_{w,d+1}^{\text{day}}=1$ indicates that the anticipated-affect question is delivered within the morning check-in.
            \item \textbf{Question:} should we treat $M^{Y,AA}_{w,d+1,1}$ and $M^{Y,AA}_{w,d+1,2}$ as the mediators?
        \end{itemize}
    \end{enumerate}
\end{frame}

\begin{frame}{Constructing state variables from mediator history}
    Each state variable is an exponentially weighted average of the corresponding mediator observed before the current decision time $(w,d,t)$, with decay factor $r=\tfrac{h-1}{h}$ and horizon $h = 2(d-1) + t-1$ (unless stated otherwise).
    \begin{enumerate}
        \item $\overline{M}^{Y,AA}_{w,d,t}=\displaystyle\sum_{i=1}^{h} r^{\,i-1}\, M^{Y,AA}_{w,i}\cdot M^{E,PJ}_{w,i} \cdot I_{w,i}^{\text{day}}$.
        \item $\overline{M}^{Y,SC}_{w,d,t}=\displaystyle\sum_{i=1}^{h} r^{\,i-1}\, M^{Y,SC}_{w,i}$.
        \item $\overline{M}^{E,PV}_{w,d,t}=\displaystyle\sum_{i=1}^{h} r^{\,i-1}\, M^{E,PV}_{w,i}$.
        \item $\overline{M}^{E,PJ}_{w,d,t}=\displaystyle\sum_{i=1}^{h} r^{\,i-1}\, M^{E,PJ}_{w,i}$.
        \item $\overline{M}^{E,FW}_{w,d,t}=\displaystyle\sum_{i=1}^{h} r^{\,i-1}\, M^{E,FW}_{w,i}$, with $h=d-1$.
    \end{enumerate}
\end{frame}

\begin{frame}{RL state}
    \begin{align}
\phi(S_{w,d,t}^{\dagger},a)
&=
\left[
1,\,
d_n,\,
t_n,\,
\widehat e_w,\,
d_n\widehat e_w,\,
t_n\widehat e_w,\,
\widehat b_w,\,
d_n\widehat b_w,\,
t_n\widehat b_w,\,
\widetilde b_w
\right] \nonumber\\
&\quad \frown
\left[
\overline{M}^{Y,AA}_{w,d,t},\,
\overline{M}^{Y,SC}_{w,d,t},\,
\overline{M}^{E,PV}_{w,d,t},\,
\overline{M}^{E,FW}_{w,d,t},\,
\overline{M}^{E,PJ}_{w,d,t},\,
C_{w,d,t}^\top
\right] \nonumber\\
&\quad \frown
a
\left[
1,\,
d_n,\,
t_n,\,
\widehat e_w,\,
d_n\widehat e_w,\,
t_n\widehat e_w,\,
\widehat b_w,\,
d_n\widehat b_w,\,
t_n\widehat b_w,\,
\widetilde b_w,\,
C_{w,d,t}^\top
\right].
\label{eq:feature-map}
\end{align}
Since $C_{w,d,t}$ has dimension 7, the state has dimension 39.
\end{frame}

\section{Environment Variants}

\begin{frame}{Testbed variants}
\begin{itemize}
    \item \textbf{Vanilla testbed:} fit the models to the MRT data and keep all model weights unchanged.
    \item \textbf{Variant 1 (constrain the sign of the weights):}
    \begin{itemize}
        \item effect of $M^Y_{w,d,t}$ on $Y_{w+1}$ non-negative: $\max\{\beta,\,0.01\}$;
        \item effect of $M^E_{w,d,t}$ on $E_{w+1}$ non-negative: $\max\{\beta,\,0.01\}$;
        \item effect of $A_{w,d,t}$ on $M^Y_{w,d,t}$ non-negative: $\max\{\beta,\,0.01\}$;
        \item effect of $A_{w,d,t}\cdot Y_w$ on $M^Y_{w,d,t}$ non-negative: $\max\{\beta,\,0.01\}$;
        \item effect of $A_{w,d,t}$ on $M^E_{w,d,t}$ non-positive: $\min\{\beta,\,-0.01\}$;
        \item effect of $A_{w,d,t}\cdot E_w$ on $M^E_{w,d,t}$ non-positive: $\min\{\beta,\,-0.01\}$;
        \item effect of $E_w$ on $M^Y_{w,d,t}$ non-negative: $\max\{\beta,\,0.01\}$.
    \end{itemize}
    \item \textbf{Other variants:} in addition to constraining the signs as in Variant 1, tune the standardized treatment effect to low (0.15), medium (0.3), or high (0.5).
\end{itemize}
\end{frame}

\begin{frame}{Standardized treatment effect}
    \begin{itemize}
        \item Let the cumulative reward for user $i$ be $G_i^{\pi} = \sum_{w=1}^{W} Y_{i,w}^{\pi}$, with $W=36$.
        \item Let the expected treatment effect be $\Delta_i = \mathbb{E}\!\left[G_i^{\pi_i^\star}\right] - \mathbb{E}\!\left[G_i^{\pi^0}\right]$.
        \item Let the variance under the zero policy be $\sigma_i^2 = \operatorname{Var}\!\left(G_i^{\pi^0}\right)$.
    \end{itemize}
    Prior work uses slightly different definitions \citep{gao25a,xu2025reinforcement}:
    \begin{enumerate}
        \item $\mathbb{E}_i\!\left[\dfrac{\Delta_i}{\sqrt{\sigma^2_i}}\right]$;
        \item $\dfrac{\mathbb{E}_i\left[\Delta_i\right]}{\operatorname{Var}_i\!\left(\mathbb{E}\left[G_i^{\pi^0}\right]\right)}$;
        \item $\dfrac{\mathbb{E}_i\left[\Delta_i\right]}{\sqrt{\mathbb{E}_i[\sigma^2_i]}}$.
    \end{enumerate}
\end{frame}

\begin{frame}{Estimating the STE}
    \begin{itemize}
        \item Approximate the optimal policy $\pi_i^\star$ with DQN.
        \item Generate 5000 trajectories per user to train DQN.
        \item Evaluate the cumulative rewards of DQN and the zero policy using 500 replicates across the 29 users.
        \item Sources of randomness within each trajectory:
        \begin{enumerate}
            \item activity suggestions are randomized with probability 0.5;
            \item residuals are resampled rather than cycled deterministically;
            \item perceived utility has no residuals and is drawn from a normal distribution;
            \item all binary variables are drawn from Bernoulli distributions.
        \end{enumerate}
        \item Each trajectory is generated under weekly queries $I_w=1$ and $J_w=1$, so DQN observes the true CAE; no belief-state update is involved.
        \item DQN uses the linearly approximated perceived utility $\widehat e_w$ rather than the ground truth $E_w$.
        \item \textbf{Questions:} should the CAE also be latent? Should we constrain the DQN randomization probability to $[0.1,\,0.9]$?
    \end{itemize}
\end{frame}

\begin{frame}{References}
    \bibliographystyle{unsrtnat}
    \bibliography{slides/refs}
\end{frame}

\end{document}
